\documentclass[12pt,a4paper]{report}

\usepackage[utf8]{inputenc}
\usepackage[T1]{fontenc}
\usepackage[french]{babel}
\usepackage{geometry}
\usepackage{xcolor}
\usepackage{array}
\usepackage{longtable}
\usepackage{float}
\usepackage{hyperref}
\usepackage{listings}
\usepackage{caption}
\usepackage{tikz}
\usetikzlibrary{arrows.meta,positioning,shapes.geometric,fit}

\geometry{margin=2.5cm}

\hypersetup{
    colorlinks=true,
    linkcolor=blue,
    urlcolor=blue
}

\definecolor{codebg}{RGB}{248,248,248}
\definecolor{codeframe}{RGB}{210,210,210}
\definecolor{mainblue}{RGB}{28,76,145}
\definecolor{lightblue}{RGB}{231,240,252}
\definecolor{lightgreen}{RGB}{232,246,238}
\definecolor{lightorange}{RGB}{255,242,225}

\lstdefinestyle{code}{
    basicstyle=\ttfamily\small,
    backgroundcolor=\color{codebg},
    frame=single,
    rulecolor=\color{codeframe},
    breaklines=true,
    showstringspaces=false,
    captionpos=b
}

\tikzset{
    box/.style={rectangle, rounded corners, draw=mainblue, thick, fill=lightblue, align=center, minimum height=1cm, text width=3.1cm},
    server/.style={rectangle, rounded corners, draw=black, thick, fill=lightgreen, align=center, minimum height=1cm, text width=3.4cm},
    note/.style={rectangle, rounded corners, draw=black!50, fill=lightorange, align=center, minimum height=0.9cm, text width=3cm},
    arrow/.style={-{Latex[length=3mm]}, thick}
}

\newcommand{\capture}[2]{
    \begin{figure}[H]
        \centering
        \fbox{
            \begin{minipage}[c][4.2cm][c]{0.86\textwidth}
                \centering
                \textbf{Capture à insérer}\\[0.3cm]
                #1\\[0.2cm]
                \nolinkurl{#2}
            \end{minipage}
        }
        \caption{#1}
    \end{figure}
}

\title{
    \textbf{Rapport professionnel du projet ChatUser}\\
    \large Comparaison des versions RMI, XML-RPC, SOAP, REST et TCP sockets
}
\author{Réalisé par : \underline{\hspace{7cm}}}
\date{Année universitaire : \underline{\hspace{5cm}}}

\begin{document}

\maketitle
\tableofcontents
\listoffigures

\chapter*{Introduction}
\addcontentsline{toc}{chapter}{Introduction}

Ce rapport présente le projet \textbf{ChatUser}, une application de discussion en ligne développée en Java et déclinée en cinq versions distribuées : \textbf{RMI}, \textbf{XML-RPC}, \textbf{SOAP}, \textbf{REST} et \textbf{TCP sockets}.

Le besoin fonctionnel est volontairement identique dans toutes les versions : un serveur gère une salle de discussion, plusieurs clients se connectent, choisissent un pseudo, envoient des messages et reçoivent les messages des autres utilisateurs. Ce qui change d'une version à l'autre, c'est la manière dont l'information circule entre le client et le serveur.

Chaque chapitre suit le même plan afin de faciliter la comparaison :

\begin{itemize}
    \item description de la technologie ;
    \item spécificité de la version ;
    \item workflow de fonctionnement ;
    \item schéma du transfert de message ;
    \item explication simplifiée avec une analogie ;
    \item méthode de test.
\end{itemize}

Les fichiers sources complets sont présents dans les dossiers du projet. Le rapport ne recopie donc pas tout le code : il explique l'architecture, le fonctionnement et les tests de manière claire et exploitable.

\chapter{Vue générale du projet}

\section{Objectif fonctionnel}

L'application doit permettre à plusieurs utilisateurs de discuter simultanément. Un utilisateur lance un client graphique, saisit un pseudo, rejoint la salle de discussion, envoie des messages et voit les messages publiés par les autres utilisateurs.

\section{Architecture commune}

Toutes les versions reposent sur une architecture client-serveur. Le client contient l'interface Swing. Le serveur contient la logique de la salle de discussion : liste des utilisateurs, réception des messages et diffusion ou mise à disposition des messages.

\begin{figure}[H]
\centering
\begin{tikzpicture}[node distance=1.3cm]
    \node[box] (c1) {Client Swing\\Alice};
    \node[box, right=1.2cm of c1] (c2) {Client Swing\\Bob};
    \node[box, right=1.2cm of c2] (c3) {Client Swing\\Charlie};
    \node[server, below=1.4cm of c2] (srv) {Serveur ChatRoom\\gestion de la salle};
    \node[note, below=1.2cm of srv] (state) {Utilisateurs\\messages\\état du chat};
    \draw[arrow] (c1) -- (srv);
    \draw[arrow] (c2) -- (srv);
    \draw[arrow] (c3) -- (srv);
    \draw[arrow] (srv) -- (state);
\end{tikzpicture}
\caption{Architecture logique commune aux cinq versions}
\end{figure}

\section{Dossiers du projet}

\begin{longtable}{|p{4cm}|p{9cm}|}
\hline
\textbf{Dossier} & \textbf{Contenu} \\
\hline
\texttt{TP0} & Version RMI. \\
\hline
\texttt{TP0-XMLRPC} & Version XML-RPC avec bibliothèque Apache XML-RPC. \\
\hline
\texttt{TP0-SOAP} & Version SOAP utilisant HTTP et des enveloppes XML. \\
\hline
\texttt{TP0-REST} & Version REST utilisant HTTP et JSON. \\
\hline
\texttt{TP0-TCP} & Version TCP sockets utilisant \texttt{ServerSocket} et \texttt{Socket}. \\
\hline
\end{longtable}

\chapter{Version RMI}

\section{Description}

RMI signifie \textit{Remote Method Invocation}. Cette technologie permet à un programme Java d'appeler les méthodes d'un objet situé dans une autre machine virtuelle Java. Dans cette version, la salle de discussion est un objet distant publié par le serveur.

\section{Spécificité de cette version}

La version RMI est la plus orientée objet. Le client ne construit pas manuellement des requêtes réseau : il récupère une référence distante vers \texttt{ChatRoom} et appelle directement ses méthodes. De plus, le client expose lui aussi un objet distant \texttt{ChatUser}. Cela permet au serveur d'appeler la méthode \texttt{displayMessage} du client pour lui pousser les messages.

\section{Workflow}

\begin{enumerate}
    \item Le serveur crée une instance de \texttt{ChatRoomImpl}.
    \item Il démarre le registre RMI sur le port \texttt{1099}.
    \item Il publie l'objet distant sous le nom \texttt{ChatRoom}.
    \item Le client récupère cet objet avec \texttt{Naming.lookup}.
    \item Le client crée son objet distant \texttt{ChatUserImpl}.
    \item Le client s'abonne à la salle avec son pseudo.
    \item Quand un message est envoyé, le serveur appelle \texttt{displayMessage} sur tous les clients.
\end{enumerate}

\section{Schéma du transfert de message}

\begin{figure}[H]
\centering
\begin{tikzpicture}[node distance=1.2cm]
    \node[box] (alice) {Client Alice\\objet ChatUser};
    \node[server, right=2.2cm of alice] (room) {Objet distant\\ChatRoom};
    \node[box, right=2.2cm of room] (bob) {Client Bob\\objet ChatUser};
    \draw[arrow] (alice) -- node[above]{\small postMessage()} (room);
    \draw[arrow] (room) -- node[above]{\small displayMessage()} (bob);
    \draw[arrow] (room) to[bend right=25] node[below]{\small displayMessage()} (alice);
\end{tikzpicture}
\caption{Transfert d'un message dans la version RMI}
\end{figure}

\section{Explication simplifiée}

RMI fonctionne comme si le client avait une télécommande vers un objet situé sur le serveur. Le client appuie sur le bouton \texttt{postMessage}, et le serveur répond en appelant directement les téléphones des autres clients avec \texttt{displayMessage}. Le point important est que tout reste dans le monde Java.

\section{Tests}

\begin{lstlisting}[style=code,caption={Tester la version RMI}]
cd TP0
ant compile
ant run-server

# Dans deux autres terminaux
cd TP0
ant run-client
\end{lstlisting}

Résultat attendu : chaque client saisit un pseudo, les messages d'arrivée sont visibles, puis les messages envoyés par Alice apparaissent chez Bob et inversement.

\capture{Serveur RMI lancé}{screenshots/rmi_server.png}
\capture{Deux clients RMI en discussion}{screenshots/rmi_clients.png}

\chapter{Version XML-RPC}

\section{Description}

XML-RPC est un protocole d'appel de procédures distantes. Le client envoie une requête HTTP contenant le nom d'une méthode et ses paramètres encodés en XML. Le serveur exécute la méthode demandée et retourne une réponse XML.

\section{Spécificité de cette version}

Contrairement à RMI, le serveur ne possède pas de référence directe vers des objets clients distants. Il expose plutôt des méthodes comme \texttt{subscribe}, \texttt{postMessage} et \texttt{getMessages}. Pour recevoir les messages, le client interroge régulièrement le serveur : c'est le principe du \textit{polling}.

\section{Workflow}

\begin{enumerate}
    \item Le serveur XML-RPC démarre sur le port \texttt{8080}.
    \item Il expose un gestionnaire appelé \texttt{ChatRoom}.
    \item Le client crée un \texttt{XmlRpcClient}.
    \item Le client appelle \texttt{ChatRoom.subscribe}.
    \item Lorsqu'un message est envoyé, le client appelle \texttt{ChatRoom.postMessage}.
    \item Toutes les 500 ms, le client appelle \texttt{ChatRoom.getMessages(lastIndex)}.
    \item Le serveur renvoie uniquement les messages ajoutés depuis le dernier index connu.
\end{enumerate}

\section{Schéma du transfert de message}

\begin{figure}[H]
\centering
\begin{tikzpicture}[node distance=1.15cm]
    \node[box] (alice) {Client Alice};
    \node[server, right=2.1cm of alice] (srv) {Serveur XML-RPC\\ChatRoom};
    \node[box, right=2.1cm of srv] (bob) {Client Bob};
    \node[note, below=1.1cm of srv] (list) {Liste des messages};
    \draw[arrow] (alice) -- node[above]{\small postMessage XML} (srv);
    \draw[arrow] (srv) -- (list);
    \draw[arrow] (bob) -- node[above]{\small getMessages} (srv);
    \draw[arrow] (srv) to[bend left=20] node[below]{\small nouveaux messages} (bob);
\end{tikzpicture}
\caption{Transfert d'un message dans la version XML-RPC}
\end{figure}

\section{Explication simplifiée}

XML-RPC fonctionne comme un guichet. Le client ne reçoit pas automatiquement un appel du serveur. Il revient régulièrement au guichet et demande : « Y a-t-il de nouveaux messages depuis mon dernier passage ? ». Le serveur répond avec la liste des nouveaux messages.

\section{Tests}

\begin{lstlisting}[style=code,caption={Tester la version XML-RPC}]
cd TP0-XMLRPC
ant compile
ant run-server

# Dans deux autres terminaux
cd TP0-XMLRPC
ant run-client
\end{lstlisting}

Résultat attendu : le serveur affiche qu'il écoute sur \texttt{http://localhost:8080/xmlrpc}. Les clients se connectent, puis les messages sont récupérés automatiquement par polling.

\capture{Serveur XML-RPC lancé}{screenshots/xmlrpc_server.png}
\capture{Deux clients XML-RPC en discussion}{screenshots/xmlrpc_clients.png}

\chapter{Version SOAP}

\section{Description}

SOAP est un protocole de services web basé sur XML. Chaque requête est placée dans une enveloppe SOAP. Cette enveloppe contient le corps de la requête, c'est-à-dire l'opération à exécuter et ses paramètres.

\section{Spécificité de cette version}

La version SOAP est plus formelle que XML-RPC. Les messages suivent une structure XML précise : \texttt{Envelope}, \texttt{Body}, puis opération. Le serveur expose aussi un WSDL minimal à l'adresse \texttt{http://localhost:8081/chat?wsdl}. Le WSDL sert à décrire le service.

\section{Workflow}

\begin{enumerate}
    \item Le serveur SOAP démarre sur le port \texttt{8081}.
    \item Le client construit une enveloppe SOAP.
    \item L'enveloppe contient une opération : \texttt{subscribe}, \texttt{postMessage}, \texttt{getMessages} ou \texttt{unsubscribe}.
    \item Le serveur parse le XML, identifie l'opération et appelle la méthode Java correspondante.
    \item Le serveur renvoie une réponse SOAP.
    \item Le client lit la réponse et met à jour l'interface.
\end{enumerate}

\section{Schéma du transfert de message}

\begin{figure}[H]
\centering
\begin{tikzpicture}[node distance=1.1cm]
    \node[box] (client) {Client SOAP};
    \node[note, right=1.6cm of client] (env) {Enveloppe SOAP\\XML};
    \node[server, right=1.6cm of env] (srv) {Serveur SOAP};
    \node[note, below=1.1cm of srv] (room) {ChatRoomImpl};
    \draw[arrow] (client) -- node[above]{\small HTTP POST} (env);
    \draw[arrow] (env) -- (srv);
    \draw[arrow] (srv) -- node[right]{\small méthode Java} (room);
    \draw[arrow] (srv) to[bend left=30] node[below]{\small réponse SOAP} (client);
\end{tikzpicture}
\caption{Transfert d'un message dans la version SOAP}
\end{figure}

\section{Explication simplifiée}

SOAP ressemble à une lettre administrative. Le message n'est pas envoyé directement : il est placé dans une enveloppe officielle avec une structure précise. Le serveur ouvre l'enveloppe, lit la demande, traite l'opération, puis répond avec une autre enveloppe.

\section{Tests}

\begin{lstlisting}[style=code,caption={Tester la version SOAP}]
cd TP0-SOAP
ant compile
ant run-server

# Verifier le WSDL dans un navigateur
http://localhost:8081/chat?wsdl

# Dans deux autres terminaux
cd TP0-SOAP
ant run-client
\end{lstlisting}

Résultat attendu : le WSDL s'affiche dans le navigateur, puis les clients échangent des messages à travers des enveloppes SOAP.

\capture{Serveur SOAP lancé}{screenshots/soap_server.png}
\capture{WSDL SOAP affiché}{screenshots/soap_wsdl.png}
\capture{Deux clients SOAP en discussion}{screenshots/soap_clients.png}

\chapter{Version REST}

\section{Description}

REST est un style d'architecture qui utilise les méthodes HTTP pour manipuler des ressources. Dans cette version, le serveur expose des endpoints et les données sont échangées en JSON.

\section{Spécificité de cette version}

La version REST est la plus proche des API web modernes. Les actions du chat sont représentées par des URLs :

\begin{longtable}{|p{3cm}|p{5cm}|p{5cm}|}
\hline
\textbf{Méthode} & \textbf{Endpoint} & \textbf{Rôle} \\
\hline
\texttt{POST} & \texttt{/chat/subscribe} & Connexion d'un utilisateur. \\
\hline
\texttt{POST} & \texttt{/chat/message} & Envoi d'un message. \\
\hline
\texttt{GET} & \texttt{/chat/messages?lastIndex=0} & Récupération des messages. \\
\hline
\texttt{POST} & \texttt{/chat/unsubscribe} & Déconnexion d'un utilisateur. \\
\hline
\end{longtable}

\section{Workflow}

\begin{enumerate}
    \item Le serveur REST démarre sur le port \texttt{8082}.
    \item Le client envoie une requête HTTP \texttt{POST} pour s'abonner.
    \item Le client envoie une requête \texttt{POST} pour publier un message.
    \item Le message est stocké côté serveur.
    \item Le client interroge régulièrement \texttt{/chat/messages}.
    \item Le serveur répond avec un JSON contenant les nouveaux messages.
\end{enumerate}

\section{Schéma du transfert de message}

\begin{figure}[H]
\centering
\begin{tikzpicture}[node distance=1.1cm]
    \node[box] (alice) {Client Alice};
    \node[server, right=2cm of alice] (api) {API REST\\HTTP + JSON};
    \node[note, right=2cm of api] (store) {Messages\\en mémoire};
    \node[box, below=1.3cm of alice] (bob) {Client Bob};
    \draw[arrow] (alice) -- node[above]{\small POST /message} (api);
    \draw[arrow] (api) -- (store);
    \draw[arrow] (bob) -- node[below]{\small GET /messages} (api);
    \draw[arrow] (api) to[bend left=20] node[right]{\small JSON} (bob);
\end{tikzpicture}
\caption{Transfert d'un message dans la version REST}
\end{figure}

\section{Explication simplifiée}

REST ressemble à un tableau d'affichage accessible par adresse. Pour publier, on envoie une fiche au serveur avec \texttt{POST}. Pour lire, on demande au serveur les nouvelles fiches avec \texttt{GET}. Le format JSON rend les échanges courts et faciles à lire.

\section{Tests}

\begin{lstlisting}[style=code,caption={Tester la version REST}]
cd TP0-REST
ant compile
ant run-server

# Test manuel possible
curl -X POST http://localhost:8082/chat/subscribe \
  -H "Content-Type: application/json" \
  -d "{\"pseudo\":\"Alice\"}"

curl -X POST http://localhost:8082/chat/message \
  -H "Content-Type: application/json" \
  -d "{\"pseudo\":\"Alice\",\"message\":\"Bonjour REST\"}"

curl http://localhost:8082/chat/messages?lastIndex=0
\end{lstlisting}

Résultat attendu : les requêtes retournent \texttt{OK}, puis la liste JSON contient les messages envoyés.

\capture{Serveur REST lancé}{screenshots/rest_server.png}
\capture{Test REST avec navigateur, curl ou Postman}{screenshots/rest_test.png}
\capture{Deux clients REST en discussion}{screenshots/rest_clients.png}

\chapter{Version TCP sockets}

\section{Description}

La version TCP sockets utilise directement la couche réseau de Java. Le serveur ouvre un \texttt{ServerSocket} sur le port \texttt{8083}. Chaque client ouvre une \texttt{Socket} vers ce serveur. Les messages sont transmis sous forme de lignes de texte.

\section{Spécificité de cette version}

Cette version est plus bas niveau que les autres. Il n'y a pas de protocole applicatif standard comme SOAP ou REST. Le protocole est défini dans le projet :

\begin{enumerate}
    \item la première ligne envoyée par le client est son pseudo ;
    \item le serveur répond \texttt{OK} ou \texttt{ERROR} ;
    \item chaque ligne suivante est considérée comme un message ;
    \item le serveur diffuse la ligne aux clients connectés.
\end{enumerate}

\section{Workflow}

\begin{enumerate}
    \item Le serveur crée un \texttt{ServerSocket}.
    \item Il attend les connexions avec \texttt{accept()}.
    \item Pour chaque connexion, il crée un thread \texttt{ClientHandler}.
    \item Le thread lit le pseudo puis les messages du client.
    \item Lorsqu'un message arrive, le serveur le diffuse à tous les clients.
    \item Quand un client ferme la fenêtre, la socket est fermée et l'utilisateur est retiré.
\end{enumerate}

\section{Schéma du transfert de message}

\begin{figure}[H]
\centering
\begin{tikzpicture}[node distance=1.1cm]
    \node[box] (alice) {Client Alice\\Socket};
    \node[server, right=2cm of alice] (srv) {ServerSocket\\port 8083};
    \node[note, below=1.1cm of srv] (thread) {Thread\\ClientHandler};
    \node[box, right=2cm of srv] (bob) {Client Bob\\Socket};
    \draw[arrow] (alice) -- node[above]{\small ligne texte} (srv);
    \draw[arrow] (srv) -- (thread);
    \draw[arrow] (thread) -- node[right]{\small broadcast} (srv);
    \draw[arrow] (srv) -- node[above]{\small ligne texte} (bob);
\end{tikzpicture}
\caption{Transfert d'un message dans la version TCP sockets}
\end{figure}

\section{Explication simplifiée}

TCP sockets ressemble à une conversation téléphonique ouverte. Tant que la ligne est active, le client peut parler au serveur. Le serveur écoute chaque ligne et la répète aux autres clients. Contrairement à REST ou SOAP, il n'y a pas d'enveloppe ou d'URL : seulement une connexion et des flux de texte.

\section{Tests}

\begin{lstlisting}[style=code,caption={Tester la version TCP sockets}]
cd TP0-TCP
ant compile
ant run-server

# Dans deux autres terminaux
cd TP0-TCP
ant run-client
\end{lstlisting}

Résultat attendu : le serveur affiche qu'il écoute sur \texttt{localhost:8083}. Deux clients peuvent se connecter, saisir des pseudos différents et échanger des messages sans polling.

\capture{Serveur TCP lancé}{screenshots/tcp_server.png}
\capture{Deux clients TCP en discussion}{screenshots/tcp_clients.png}

\chapter{Comparaison générale}

\section{Tableau comparatif}

\begin{longtable}{|p{2.4cm}|p{2.3cm}|p{2.3cm}|p{2.3cm}|p{2.3cm}|p{2.3cm}|}
\hline
\textbf{Critère} & \textbf{RMI} & \textbf{XML-RPC} & \textbf{SOAP} & \textbf{REST} & \textbf{TCP} \\
\hline
Transport & RMI Java & HTTP & HTTP & HTTP & TCP brut \\
\hline
Format & Objets Java & XML & XML SOAP & JSON & Texte \\
\hline
Port & 1099 & 8080 & 8081 & 8082 & 8083 \\
\hline
Réception & Callback & Polling & Polling & Polling & Connexion ouverte \\
\hline
Interopérabilité & Faible & Moyenne & Bonne & Très bonne & Dépend du protocole \\
\hline
Complexité & Moyenne & Moyenne & Élevée & Faible & Moyenne \\
\hline
\end{longtable}

\section{Analyse}

RMI est très pratique dans un environnement Java pur. XML-RPC et SOAP montrent deux approches historiques de services web en XML. REST est plus léger et plus facile à tester. TCP sockets donne la meilleure compréhension des échanges réseau, mais demande de gérer soi-même le protocole.

\chapter{Guide de test complet}

\section{Vérifier l'environnement}

\begin{lstlisting}[style=code,caption={Vérification Java et Ant}]
java -version
ant -version
\end{lstlisting}

\section{Scénario de test commun}

Pour chaque version, le scénario de validation est :

\begin{enumerate}
    \item compiler le projet ;
    \item lancer le serveur ;
    \item lancer deux clients ;
    \item saisir deux pseudos différents ;
    \item envoyer un message depuis le premier client ;
    \item vérifier que le second reçoit le message ;
    \item envoyer un message depuis le second client ;
    \item vérifier que le premier reçoit le message ;
    \item fermer un client ;
    \item vérifier que le serveur continue à fonctionner.
\end{enumerate}

\section{Commandes de compilation globale}

\begin{lstlisting}[style=code,caption={Compiler toutes les versions}]
cd TP0
ant compile

cd ../TP0-XMLRPC
ant compile

cd ../TP0-SOAP
ant compile

cd ../TP0-REST
ant compile

cd ../TP0-TCP
ant compile
\end{lstlisting}

\section{Points à vérifier pendant les tests}

\begin{longtable}{|p{5cm}|p{8cm}|}
\hline
\textbf{Point testé} & \textbf{Résultat attendu} \\
\hline
Connexion d'un client & Le client demande un pseudo et rejoint la salle. \\
\hline
Connexion de deux clients & Chaque client voit l'arrivée de l'autre. \\
\hline
Envoi d'un message & Le message est visible chez les autres clients. \\
\hline
Message local & Le client émetteur voit son propre message sous une forme claire, par exemple \texttt{Vous:}. \\
\hline
Déconnexion & La fermeture d'un client ne coupe pas le serveur. \\
\hline
Pseudo déjà utilisé & La version doit refuser ou gérer correctement le doublon selon son implémentation. \\
\hline
Serveur arrêté & Le client doit afficher une erreur de connexion. \\
\hline
\end{longtable}

\chapter*{Conclusion}
\addcontentsline{toc}{chapter}{Conclusion}

Le projet ChatUser montre qu'une même application peut être réalisée avec plusieurs modèles de communication. RMI donne une vision objet distante, XML-RPC et SOAP illustrent les appels distants en XML, REST représente une approche web moderne, et TCP sockets expose les mécanismes réseau fondamentaux.

La comparaison met en évidence que chaque technologie répond à un besoin différent. Pour un projet Java pur, RMI est naturel. Pour une API simple et testable, REST est plus adapté. Pour un cadre formel de service web, SOAP reste intéressant. Pour comprendre les bases du réseau, TCP sockets est la version la plus pédagogique.

\end{document}
